\section{The finite algebra at the marked point}
\label{sec:marked-hosts}

The support comparison in Section~\ref{sec:linked-components} took us to an abelian CM
extension \(L\). We now need a point there that will contribute to a length computed
over the original coefficient DVR \(\mathcal O\).
Two base changes intervene. The residual curve
points have to descend from the enlarged frame, and a point of
the graded algebra in characteristic zero has to give an integral
representation to which the automorphy comparison applies.
The marked point algebra gives the integral model used for the
descent and for the length bound in the next two appendices. The separated frame of Proposition~\ref{prop:separated-marked-host} carries the mark \eqref{eq:tag-h1}. Its degree-zero algebra over the original coefficient DVR is the object of Proposition~\ref{prop:marked-invariant-algebra}.

\paragraph{The fields and the mark.}
The field and its mark are specified by characters. With
\(e\ge1\) and \(\psi\) of even modulus \(N\) satisfying \(\psi(-1)=(-1)^e\), form the characters
\[
 \Psi=\epsilon^e\psi,\qquad P=\bar\Psi,\qquad \Psi(c)=1 .
\] Here \(c\) denotes complex conjugation.
Let \(\eta\) cut out the imaginary quadratic field \(K\). The
linking cutoff is \(\ell N|D_K|\) in the general-data case and
\(\ell N\) for \(K=\Q(i)\). Indeed an inertia group on which
\(\eta\) is nontrivial lies over a divisor of \(D_K\), and the
Gaussian character has conductor four. Let \(B\) be the chosen
place over \(\ell\). The marking uses
\[
       g_*\in D_B\cap G_L,\qquad P(g_*)^2\ne1 .
       \tag{H1}\label{eq:tag-h1}
\]
The level of the final packet will be read from the inertia of
the original representation.

\begin{proposition}[A separated marked frame]\label{prop:separated-marked-host}
Under the general-data or Gaussian host hypotheses, the
construction yields a totally real abelian \(L^+\) and its
CM extension \(L=L^+K\). The extension \(L/L^+\) splits
at the places above the primes dividing \(M\) and is unramified
at every finite place. At every place \(v\mid\ell\) of \(L^+\),
\[
                     [L_v^+:\mathbb Q_\ell]>3C_0+7.
\]
The distinguished mark satisfies
\textup{\eqref{eq:tag-h1}}. Both signed residual extension classes with
characters \(\eta P^{-a}\), \(a=1,2\), restrict nontrivially
to the marked frame with its prescribed inertia conditions.
\end{proposition}

\begin{proof}
\proofstep{Matching the quadratic characters locally.}
Here are the choices in the construction of the frame underlying
\eqref{eq:tag-h1}. Write \(K=\Q(\sqrt{-D})\), \(D>0\), and denote the linking
cutoff by \(M\). Choose a prime \(p>4D+M\) congruent to \(-1\)
modulo \(4M\). The real quadratic field
\(\Q(\sqrt{Dp})\) has the same local character as \(K\) on
decomposition groups above \(M\). Their defining square classes have ratio \(Dp/(-D)=-p\). For an odd prime \(v\mid M\), the congruence \(-p\equiv1\pmod v\) lifts the simple root \(1\) of \(X^2+p\) to \(\mathbb Z_v\). If \(2\mid M\), then \(8\mid4M\). Put \(u=-p\equiv1\pmod8\) and start with \(x_3=1\).

For \(n\ge3\),
\[
 (x_n+e2^{n-1})^2\equiv x_n^2+e2^n x_n\pmod {2^{n+1}} .
\]
If \(x_n^2\equiv u\pmod {2^n}\), exactly one of \(e=0,1\) corrects the next bit because \(x_n\) is odd. The resulting sequence converges to a square root in \(\mathbb Z_2^\times\).
This real quadratic field also ramifies at \(p\), where \(K\)
is unramified. Denote its fixing group by \(H_0\). The general
frame data include an abelian quotient of \(\ell\)-power order
with kernel \(H_b\) on which the away-from-\(\ell\) inertia
character has the same kernel as its reduction. On the Gaussian
branch \(H_b\) kills the \(\ell\)-power part of the character
after removing the factor at the place \(\ell\).

On inertia at \(\ell\), the condition \(P=\bar\epsilon^j\) determines the reduction of the finite factor. Use the Teichm\"uller lift
\[
 \chi'=[P\bar\epsilon^{-j}],\qquad\bar\epsilon^j\bar\chi'=P .
\]
The finite characters \(\chi'\) and \(\eta\) factor through
roots of unity of order dividing some \(M^a\ell\). Both are
unramified at \(\ell\), so its \(\ell\)-power factor can be
removed from that modulus by cyclotomic inertia. Choose a
second prime \(w\) congruent to one modulo both the resulting
modulus and \(M^{a+1}\ell\), outside the divisors of those
integers. The degree-\(\ell\) subfield of \(\Q(\mu_w)\) gives
the order-\(\ell\) character \(\theta\) of the seed construction.

For \(m\ge1\), let \(H_w\) fix that subfield, and \(H_m\) fix the
degree-\(\ell^m\) layer in \(\Q(\mu_{\ell^{m+1}})\). Put
\[
       U_m=H_0\cap H_w\cap H_b\cap H_m,\qquad
       G_L=U_m\cap\ker\eta .
       \tag{H2}\label{eq:tag-h2}
\]
All quotients defining \(U_m\) are abelian and complex conjugation
lies in each kernel. Thus \(U_m=G_{L^+}\), \(L=L^+K\). The
local comparison for \(H_0\) gives splitting of \(L/L^+\) over
\(M\). Outside \(M\), \(\eta\) is unramified, so \(L/L^+\) is
everywhere unramified at finite places. The layer at \(w\)
kills \(\theta\) on \(G_L\).

\proofstep{Selecting the real abelian layers.}
For the order of the choices in \eqref{eq:tag-h2}, let \(E\) be fixed by
\[
       H_0\cap H_w\cap H_b\cap\ker\eta\cap\ker P
                    \cap\ker\bar\epsilon .
\]
It is abelian over \(\Q\) and contains \(\mu_\ell\). On adjoining
the \(m\)-th layer the corresponding intersection fixes
\(E(\mu_{\ell^{m+1}})\). The bounded \(\ell\)-rank of the class
groups in this tower follows from Ferrero--Washington for the
cyclotomic tower of an abelian field and the Iwasawa module
structure theorem \cite[Chapter~13]{washington-cyclo}. Let \(E_\infty\) be their union and \(X\) its unramified abelian pro-\(\ell\) Galois module. The vanishing of \(\mu\) makes \(X\) finitely generated
over \(\mathbb Z_\ell\). Above a sufficiently deep layer \(E_m\),
one of the primes over \(\ell\) is totally ramified in
\(E_\infty/E_m\), since cyclotomic inertia has open image.

Write \(H_m^{\rm ur}\) for the unramified abelian
\(\ell\)-class field of \(E_m\). Every nontrivial finite intermediate extension in this part of the tower is ramified at the totally ramified prime over \(\ell\). Since \(H_m^{\rm ur}\) is unramified there,
\[
 H_m^{\rm ur}\cap E_\infty=E_m .
\] Unramifiedness survives base
change, so restriction gives
\[
 X\twoheadrightarrow\operatorname{Gal}
            (H_m^{\rm ur}E_\infty/E_\infty)
      \simeq\operatorname{Gal}(H_m^{\rm ur}/E_m).
\]
The \(\ell\)-rank is bounded by the number of
\(\mathbb Z_\ell\)-generators of \(X\). The finite
set of earlier layers can be included in the same bound.
By unramified class field theory this bounds
the dimension of the unramified mod-\(\ell\) character space by
a single \(C_0\) for all \(m\). Only now take
\(\ell^m>3C_0+7\). Inertia at \(\ell\) surjects onto the
degree-\(\ell^m\) layer, so every relevant local degree is at
least \(\ell^m\).

To construct the element in \eqref{eq:tag-h1}, choose \(x\in D_B\cap\ker\eta\) with \(P(x)^2\ne1\) as in the local separation hypothesis.
It fixes \(\sqrt{Dp}\) by the local comparison. Raise \(x\)
to a common \(\ell\)-power killing the three remaining finite
quotients in \(U_m\). Write \(d=\operatorname{ord}(P(x)^2)>1\). Since \(\gcd(d,\ell)=1\),
\[
 \operatorname{ord}\bigl(P(x)^{2\ell^a}\bigr)=d.
\]
Thus the resulting power \(x_*=x^{\ell^a}\) satisfies
\[
 x_*\in D_B\cap G_L,\qquad
 P(x_*)^2=(P(x)^2)^{\ell^a}\ne1 .
\]
The second equality uses the prime-to-\(\ell\) order of \(P(x)^2\),
on whose cyclic subgroup the indicated power is an automorphism.

\proofstep{Local separation by inertia and Frobenius in the general and Gaussian cases.}
Here are the two separation tests used above. For the General case of the rank-three trace estimate, inertia at the ramified prime of \(K\) away from
\(\ell N\) gives an element \(s\) with
\(P(s)=\bar\epsilon(s)=1\), \(\eta(s)=-1\). On inertia at \(\ell\) the residual finite factor has square one and \(\eta=1\). Thus here
\[
 P^2=\bar\epsilon^{2e}\ne1,\qquad \ell-1\nmid 2e .
\]

On the Gaussian branch let \(2^c\) be the
conductor of the dyadic component \(\psi_2\). Here \(c\ge4\).
Its image on units congruent to one modulo four is a nontrivial
cyclic 2-group. Some such unit maps to \(-1\). Multiplying
\(-1\) by that unit if \(\psi_2(-1)=-1\) gives a unit
congruent to \(-1\) modulo four in the kernel of \(\psi_2\).
Its image under dyadic reciprocity gives the required \(s\).
In either case \(\chi'(s)=1\) and \(\eta(s)=-1\).

For the other Gaussian test, the base data give an integer
\(o\ge1\) such that the odd-prime components of the
prime-to-\(\ell\) part of \(\psi\) have order dividing \(2o\).
For each exceptional prime and for \(p=\ell\) the dyadic conductor bound reads
\[
 8op^2<2^c .
\] These primes are odd. Choose a
Frobenius lift \(h\) acting by \(p\) on prime-to-\(p\) roots
and with \(\bar\epsilon(h)=1\). At \(p=\ell\) use cyclotomic
inertia to arrange the latter condition. At the other primes
it follows from \(p\equiv1\pmod\ell\). If \(P(h)^4=1\),
reduction is faithful on prime-to-\(\ell\) roots of unity,
so raising to the power \(4o\) gives
\(\psi_2(p)^{4o}=1\).

Let \(s_2\) and \(d\) be the two valuations that determine the effect of raising the odd prime \(p\) to successively doubled exponents:
\[
 s_2=v_2(p^2-1),\qquad d=v_2(o).
\]
For \(a\in1+4\mathbb Z_2\) one has \(v_2(a+1)=1\), so
\[
 v_2(a^2-1)=v_2(a-1)+v_2(a+1)=v_2(a-1)+1 .
\]
The geometric sum for an odd power is a unit. Starting
from \(a=p^2\) and using \(4o=2\cdot2^{d+1}(o/2^d)\)
therefore gives
\[
               v_2(p^{4o}-1)=v_2(p^2-1)+1+d .
\]
The same order formula in the finite dyadic quotients
shows that \(p^{4o}\) generates the subgroup congruent to
one modulo \(2^{s_2+1+d}\) in each such quotient. This forces
\[
 2^c\le2^{s_2+1+d}\le 2o(p^2-1),
\] a contradiction.
Since \(x=h^2\), multiplicativity of all three characters gives
\[
 \eta(x)=1,\qquad P(x)^2=P(h)^4\ne1,\qquad\bar\epsilon(x)=1 .
\] At \(p=\ell\) this is the starting
element for \eqref{eq:tag-h1}.

\proofstep{The signed residual cocycles on the marked frame.}
One use of these tests is the construction of the two
residual extension classes on the marked frame. For
\(a=1,2\) put \(\kappa=\eta P^{-a}\). This character is
odd, and differs from \(\bar\epsilon\) on \(s\) above.
It is nontrivial on \(U_m\cap D_B\) by \eqref{eq:tag-h1}. Let \(S\)
consist of \(\ell,w\) and the ramified primes of \(\kappa\).
For the Selmer group of \(k(\kappa)\), the local subspace in \(H^1(\mathbb Q_v,k(\kappa))\) is 
\[
 L_v=\begin{cases}H^1(\mathbb Q_v,k(\kappa)),&v=w,\\ 0,&v\in S\setminus\{w\},\\ H^1_{\rm unr}(\mathbb Q_v,k(\kappa)),&v\notin S.\end{cases}
\]
 The dimension difference
with the dual Selmer group is one by the Greenberg--Wiles
formula. Indeed both global invariant spaces vanish, and
the real invariant space vanishes. At \(w\) both \(\kappa\) and \(\bar\epsilon\) are trivial. The local contribution of the full condition is therefore
\[
 \dim L_w-\dim H^0_w=2-1=1 .
\]

At every other prime of \(S\) the
invariant space is zero. A nonzero resulting class remains
nonzero on \(U_m\) by inflation--restriction, since the
coefficient module has no \(U_m\)-invariants. On inertia
away from \(\ell,w\) in \(\ker\eta\cap\ker P\) its cocycle
is zero, either by the unramified condition or by local
triviality and \(\kappa=1\) on that inertia. At \(w\) the
restriction to inertia is additive. Its \(\ell\)-primary
tame quotient is procyclic and the degree-\(\ell\) field
fixed by \(H_w\) is totally ramified there. Thus it vanishes
on the inertia in \(U_m\) at \(w\) as well. Its restriction
at \(B\) is a coboundary. These are exactly the local
conditions needed to descend the two classes to the frame
quotient.

\end{proof}

\paragraph{Inertia and the abelian quotients of the fitted frame.}
The generator count uses these inertia quotients.

Let \(G\) be the quotient of \(U_m=G_{L^+}\) by the closed
normal subgroup generated by the away-from-\(\ell\) inertia
elements on which both \(\eta\) and \(P\) are one. Its positive
component is the image of \(G_L\). On the away-from-\(\ell\)
inertia of \(U_m\), the full character \(\Psi\) has prime-to-
\(\ell\) order, since \(H_b\) kills the other part and the
cyclotomic character is unramified there. Reduction on that
prime-to-\(\ell\) subgroup is injective. Thus the geometric
inertia condition on \(G\) kills the same kernel as the
ordered residual character at these places.

For each place of \(L^+\) above \(\ell\), use a representative
conjugate to \(B\) by an element of \(\ker\eta\). Let \(g\) be any conjugating element. If its sign is negative, multiply on the left by complex conjugation \(c\in U_m\):
\[
 \eta(g)=-1\quad\Longrightarrow\quad \eta(cg)=(-1)(-1)=1 .
\]
This preserves the \(U_m\)-orbit of the representative. The
places split in \(L/L^+\), so the chosen decomposition group
of \(U_m\) is also the decomposition group on the indicated
side over \(L\). These are the local groups mapped into the
frame. They are the actual decomposition groups before the
projection to \(G\).

There is a nontrivial mod-\(\ell\) cyclotomic value in each
of these local groups. On rational inertia at \(\ell\), the
quadratic character \(\eta\) is trivial in both host cases,
and so is the character of \(H_0\) by the local comparison.
Start with a nontrivial mod-\(\ell\) cyclotomic value and take
an \(\ell\)-power to land in \(H_b\cap H_w\cap H_m\).
Since the nontrivial residue value \(a\) lies in \(\mathbb F_\ell^\times\), raising it to an \(\ell\)-power preserves its exact order:
\[
 \operatorname{ord}(a^{\ell^t})=\operatorname{ord}(a)>1 .
\] Conjugating supplies
the element at every selected place.

In particular the
completions on that side contain no primitive \(\ell\)-th
root of unity. There is also a global element of \(G_L\)
on which \(P=1\) but \(\bar\epsilon\ne1\). Take the element
\(s\) of the preceding separation test, with
\((\eta,P,\bar\epsilon)(s)=(-1,1,1)\). At the product of \(s\) with complex conjugation the three character values are \((1,1,-1)\).
If needed, correct the character of \(H_0\) by an inertia
element at its new ramified prime \(p\). All three characters are trivial on this inertia group. An \(\ell\)-power again
lands in the remaining three kernels. This proves the
global separation required at the auxiliary level.

\begin{lemma}[The fitted unramified character bounds]
\label{lem:fitted-abelian-quotient}
In the frame of Proposition~\ref{prop:separated-marked-host}, let a
finite abelian \(\ell\)-group character of the positive component
vanish on \(P=1\) inertia at \(\ell\) and all its conjugates, including
complex conjugates. Its restriction \(\chi|_{G_L}\) factors as 
\[
 G_L\longrightarrow\operatorname{Cl}(L)[\ell^\infty]\longrightarrow\operatorname{im}\chi.
\]
 On the Gaussian sign kernel
with vanishing on the marked inertia intersections the unramified
pullback is trivial. For residual coefficients formed from
\(\eta,P,\bar\epsilon\), a class whose restrictions kill inertia
over \(E_m=E(\mu_{\ell^{m+1}})\) restricts to a character of
\(\operatorname{Cl}(E_m)\).
\end{lemma}

\begin{proof}
\proofstep{The unramified quotient and the class number.}
Consider a continuous finite abelian \(\ell\)-group character on
the positive component of \(G\). Suppose it vanishes on the
indicated \(P=1\) inertia at \(\ell\) and its conjugates,
including the conjugates by complex conjugation. The quotient
of inertia by that kernel has order dividing \(\#k-1\), where
\(k\) is the residue field. This quotient has trivial image
in the \(\ell\)-group. At every place away from \(\ell\),
the inertia of \(G_L\) already has image of order dividing
\(\#k-1\), by the definition of \(G\). The resulting character on \(G_L\) is unramified with \(\ell\)-primary image.

By class field theory its factorization at this stage goes through
\[
 G_L\longrightarrow\operatorname{Cl}(L)[\ell^\infty]\longrightarrow\operatorname{im}(\chi),
\] Its exponent gives one
positive annihilating integer for every truncation of a
principal-unit character. An involution-invariant character
with the specified vanishing on inertia on one side has
it on the other side as well. This is the case for the
middle character in the exceptional corner. On the original
Gaussian quotient \(G_{\Q,S}\), use instead the sign kernel
giving \(G_{\Q(i)}\). A character trivial on the conjugates
of all the indicated intersections with inertia pulls back
to an everywhere unramified character of \(\Q(i)\). The norm-Euclidean property of the Gaussian integers makes every ideal principal. Hence
\[
 \operatorname{Cl}(\mathbb Q(i))=1,\qquad h_{\mathbb Q(i)}=1 .
\]
Here the pullback character is trivial at each finite
truncation. \cite{ray-class}\label{cor:ordinary-fitted-arithmetic-1}

\proofstep{Restriction of residual classes.}
The field \(E_m=E(\mu_{\ell^{m+1}})\) is
obtained from \(L^+\) by adjoining the fields of
\(\eta,P,\bar\epsilon\), of degrees prime to \(\ell\).
For any residue character formed from these three, restriction
of \(H^1\) to \(G_{E_m}\) is injective by
inflation--restriction and prime-to-\(\ell\) averaging.
There the coefficients are constant. In particular when the
local conditions kill inertia on this subgroup, the classes
embed in the additive character group of the class group of
\(E_m\). 
\end{proof}

\paragraph{Uniform generator bounds and local deformation data.}
With the abelian quotients fixed, we can count generators and record the local
conditions for the relative cohomology calculation.

By Lemma~\ref{lem:fitted-abelian-quotient}, the strict residual
classes over \(E_m\) embed in the additive character group of its
class group. Its dimension over \(k\) is bounded by \(C_0\).
Apply this to the strict group for \(\eta P^{-2}\) over
\(L^+\) and to that for \(P^{-1}\) over \(L\). In the
first case impose a coboundary at every marked local
group.

In the second impose it at both conjugate places.
The local character is nontrivial in each case by
\(P(g_*)^2\ne1\), transported to the marked local groups.
Indeed, if \(b,b'\) trivialize the same local coboundary for the nontrivial character \(\chi\), then for an element \(g\) with \(\chi(g)\ne1\),
\[
 (\chi(g)-1)(b-b')=0\quad\Longrightarrow\quad b=b' .
\]
Keeping these scalars in the relative complex adds no
dimension to its first cohomology. After restriction to
\(G_{E_m}\), a coboundary is zero. At the other finite
places inertia below \(G_{E_m}\) is killed by the frame
quotient. The same two bounds by \(C_0\) therefore hold for
the relative strict groups.

\begin{lemma}[Additive character bounds and deformation directions]
\label{lem:fitted-local-generators}
Use the fitted frame quotient of Lemma~\ref{lem:fitted-abelian-quotient}.
For \(\eta P^{-2}\) over \(L^+\) require a coboundary at every marked
local group, and for \(P^{-1}\) over \(L\) require it at both conjugate
places. Impose unramified conditions at the remaining finite places.
Writing \(S_{L^+}(\eta P^{-2})\) and \(S_L(P^{-1})\) for the resulting
relative strict character spaces, one has
\[
  \max\bigl(\dim S_{L^+}(\eta P^{-2}),\ \dim S_L(P^{-1})\bigr)\le C_0.
\]
For the trivial character use the full local conditions at \(\ell\)
and the unramified conditions elsewhere. The ordinary local deformation
directions can be chosen independently at the marked places.
\end{lemma}

\begin{proof}
\proofstep{The bound for additive characters.}
For the trivial character, every additive character on
\(G\) is unramified away from \(\ell\). Take the Selmer group
over \(L^+\) with full conditions at \(\ell\) and unramified
conditions elsewhere. The dual coefficient line is
\(k(1)\). Its global invariant space is zero by the
cyclotomic test. Its local invariant spaces at \(\ell\) are
also zero. The Greenberg--Wiles formula gives the difference
of the two Selmer dimensions as
\[
           1+\sum_{v\mid\ell} d_v-[L^+:\Q]=1 .
\]
Indeed the local full condition contributes \(d_v\) by
Tate duality and the local Euler characteristic. An
unramified finite-place condition contributes zero.

At each real place the invariant dimension in the trivial
line is one. The dual Selmer group has zero local condition
at \(\ell\) and is unramified at the remaining finite
places, by orthogonality of unramified conditions away
from \(\ell\). Restriction to \(G_{E_m}\) bounds it by
\(C_0\). This gives the \(1+C_0\) bound for the first
cohomology with trivial coefficients. All the bounds use
the same frame as the earlier locally split cocycles.
The class rank bound precedes the choice of the deep
layer. If \(g\) is the number of these local places, summing the degree lower bounds gives
\[
 g\ell^m\le\sum_{v\mid\ell}[L^+_v:\mathbb Q_\ell]=[L^+:\mathbb Q]=d .
\]
\cite{ferrero-washington,washington-cyclo,milne-duality}

The Hecke levels use a second group, the Galois group
unramified outside the level set over \(L^+\), and its
quadratic kernel. They map to the fitted frame since the
closed inertia subgroup outside those levels is in the
kernel of the frame projection. Both sign and cyclotomic
character descend to the level group. For the determinant and patching constructions the finiteness condition on each level group and local group \(G\) is 
\[
 \dim_{\mathbb F_\ell}\operatorname{Hom}_{\rm cont}(U,\mathbb F_\ell)<\infty\qquad(U\subseteq G\text{ open}).
\]
 To see this for an open subgroup \(U\)
of a level group \(\Gamma\), apply global cohomological
finiteness to
\(\operatorname{Coind}_{U}^{\Gamma}\mathbb F_\ell\).
This is the finite continuous permutation module on
\(\Gamma/U\) (equivalently on the opposite cosets for the
right action). The level set is finite and includes the places
above \(\ell\). Shapiro gives
\[
 \operatorname{Hom}_{\rm cts}(U,\mathbb F_\ell)
   =H^1(U,\mathbb F_\ell)
   \simeq H^1(\Gamma,\operatorname{Coind}_{U}^{\Gamma}
                                    \mathbb F_\ell).
\]
The last group is finite by the global finiteness theorem.
This applies within the open quadratic kernel as well.

For a
local decomposition group use the same formula and local
finiteness for its completion. It also gives the finiteness
on the original rational level for the Gaussian base. If \(P_U\)
is the maximal pro-\(\ell\) quotient, these homomorphisms are the
dual of \(P_U/\overline{P_U^\ell[P_U,P_U]}\). The Frattini
criterion now makes \(P_U\) topologically finitely generated. This is the finite generation used in
the two universal determinant constructions.
\cite{milne-duality,chenevier}\label{cor:ordinary-fitted-arithmetic-2}

\proofstep{Independent deformation directions above the coefficient prime.}
At a selected \(\ell\)-place the maximal pro-\(\ell\) quotient
of the abelianized inertia image is, by local reciprocity,
the principal units of \(L_w\). At sufficiently large depth \(s\), logarithm identifies the deep units with an additive lattice of finite index in the integers:
\[
 \log(1+\mathfrak p_w^s)=\mathfrak p_w^s,\qquad \operatorname{rank}_{\mathbb Z_\ell}\mathfrak p_w^s=[L_w:\mathbb Q_\ell].
\]
The torsion is exactly the group of \(\ell\)-power roots of
unity, which is trivial by the cyclotomic test above.
Inside inertia the intersection of the kernels of \(P\) and
\(\bar\epsilon\) has prime-to-\(\ell\) index. Its image in
the free abelian pro-\(\ell\) unit group is the whole group,
since the cokernel has order dividing that index. On this
intersection the remaining cyclotomic map is the map
from the unit group to \(1+\ell\mathbb Z_\ell\). The image is
nonzero and open, by finite index in rational inertia.

If \(U_v\) denotes this pro-\(\ell\) unit lattice,
the remaining cyclotomic character gives an exact sequence
\[
 0\longrightarrow \ker\epsilon_\ell|_{U_v}
  \longrightarrow U_v\xrightarrow{\epsilon_\ell}
                   \epsilon_\ell(U_v)\longrightarrow1 .
\]
The image is free of rank one over \(\mathbb Z_\ell\),
so it splits by lifting one generator. The kernel is
therefore free of rank \(d_v-1\). Choose its basis along
with that one lift. Lift all \(d_v\) basis elements to the
intersection in inertia. Each has \(P=1\), and all but the
distinguished lift have full cyclotomic value one.
Vanishing on this basis makes a finite abelian
\(\ell\)-group character of the local group vanish on
inertia. This is the generator choice used to compare
two locally coincident labels on
\(N_v=I_v\cap\ker\epsilon\).
On the indicated inertia group \(P=\bar\epsilon^j\), so the defining kernel of \(N_v\) gives
\[
 P|_{N_v}=\bar\epsilon^j|_{N_v}=1 .
\]

For the decomposition group there is one further unramified
generator, obtained by lifting a uniformizer. For a
prescribed determinant weight at rank \(n\), divide each
inertia character by the Teichm\"uller lift of its residue.
Maps from the resulting free lattice to principal units
are specified by the values of a basis, which can be any
\(1+X_{vi}\) in the target. There are \(d_v\) free inertia coordinates for each of the \(n\) ordered characters at \(v\). With \(d=\sum_{v\mid\ell}d_v\), the weight algebra has the presentation 
\[
 \mathcal O[[X_1,\ldots,X_{nd}]],\qquad \sum_{v\mid\ell}n d_v=nd.
\]
 At a finite place away from \(\ell\) the
pro-\(\ell\) tame inertia quotient is procyclic. A chosen
lift of Frobenius together with inertia generates the
decomposition group topologically, since the unramified
quotient is generated densely by Frobenius. These statements
use inertia in the selected decomposition group, with
conjugacy transporting the generators to another choice of
place. \cite{ray-class,milne-duality}\label{cor:ordinary-fitted-arithmetic-3}

\end{proof}

\paragraph{The relative Euler inequality on the marked quotient.}
The generator bounds alone do not compare deformations with their obstructions. Here
the ordinary flags enter through the local quotients in a relative Euler
characteristic.

\begin{proposition}[Relative Euler inequality on the fitted quotient]
\label{prop:fitted-relative-euler}
Let \(\ell\) be odd. Let \(G\) be the quotient of \(G_{L^+}\) by
the closed normal subgroup generated by \(I_q\cap\ker\eta\cap\ker P\)
for \(q\nmid\ell\), and
let \(W\) be a finite-dimensional continuous representation of \(G\)
over a finite field \(k\) of characteristic \(\ell\). For each
\(v\mid\ell\) in \(L^+\), set \(d_v=[L_v^+:\mathbb Q_\ell]\)
and let \(b_v\subset W\) be stable under its marked group.
Assume \(\operatorname{Hom}_{G}(W,k(1))=0\).
Put \(Q_v=W/b_v\), and let \(H_C^i\) be the cohomology of the
cone of restriction from global cochains to local cochains
of \(Q_v\), shifted by \(-1\). In degree one this keeps a
global cocycle and local trivializations modulo \(b_v\).
If \(c_w\) denotes conjugation at a real place, the desired
bound is
\begin{equation}\label{eq:fitted-relative-euler}
 \dim H_C^2+\dim H_C^0+\sum_{v\mid\ell}d_v\dim b_v
 \ \le\ \dim H_C^1+\sum_{w\mid\infty}\dim W^{c_w}.
\end{equation}

\end{proposition}

\begin{proof}
\proofstep{Comparison with the arithmetic cohomology groups.}
To compare with global duality, first enlarge the set of
allowed ramification of \(L^+\) from the \(\ell\)-places to a
finite set \(S'\) including the ramification of \(\eta\) and
\(P\). Denote the unrestricted group on \(S'\) by \(\widetilde G\).
Its map to \(G\) has kernel \(J\) normally generated by the
images of \(J_q=I_q\cap\ker\eta\cap\ker P\) at
\(q\in S'\), \(q\nmid\ell\). The quotient \(I_q/J_q\) has order
prime to \(\ell\). Thus a class of \(H^2(G,W)\) has inflated
restriction zero at each such \(q\). Its local inflation
factors via the decomposition group modulo \(J_q\), an extension
of the procyclic unramified quotient by a finite group of
order prime to \(\ell\), whose \(\ell\)-cohomological dimension
is at most one.

Restriction gives an injection
\[
 \operatorname{Hom}_{\rm cont}(J,W)^{\widetilde G}
 \longrightarrow \bigoplus_{q\nmid\ell}\operatorname{Hom}_{D_q}(J_q,W).
\]
The restrictions determine a homomorphism on the closed
normal generators. The dimension of the local summand is
\(h^2(D_q,W)\), with \(h^i\) denoting dimension of continuous
cohomology. In fact wild inertia has order prime to \(\ell\)
in its finite quotients, and the pro-\(\ell\) part of tame
inertia is \(\mathbb Z_\ell(1)\), unchanged by passing to
the prime-to-\(\ell\) index subgroup. Equivariance asks for the
Frobenius-fixed space of \(W(-1)^{I_q}\). Local duality counts
the coinvariants of the same space under Frobenius, of the
same dimension. Averaging over \(I_q/J_q\) identifies inertia
invariants and coinvariants in this comparison.

\proofstep{The relative tangent and obstruction terms.}
Write \(a=h^1(\widetilde G,W)-h^1(G,W)\), and let \(T\)
denote the invariant homomorphism space on \(J\) just
considered. The relevant five-term sequence is
\[
 0\to H^1(G,W)\to H^1(\widetilde G,W)\to T
 \to H^2(G,W)\to H^2(\widetilde G,W).
\]
The image of \(H^1(\widetilde G,W)\) in \(T\) has dimension \(a\);
therefore the kernel at \(H^2(G,W)\) has dimension \(\dim T-a\).
Let \(K_2\) be the kernel of the map from \(H^2(G,W)\) to
\(\bigoplus_{v\mid\ell}H^2(D_v,W)\).
The image of \(K_2\) is locally zero at every place of
\(S'\), including the deleted primes by the inflation
argument.

In particular \(H^2(G,W)\) is finite-dimensional,
and the Poitou--Tate sequence on \(S'\) gives
\[
 \dim K_2 \le \dim T-a+h^2(\widetilde G,W)
                         -\sum_{v\in S'}h^2(D_v,W).
\]
The global dual invariant term vanishes by hypothesis.
Use the bound on \(\dim T\) and the global Euler characteristic
for \(\widetilde G\). It follows that
\[
 \dim K_2+\sum_{v\mid\ell}h^2(D_v,W)
        +h^0(G,W)+d\dim W
 \le h^1(G,W)+\sum_{w\mid\infty}\dim W^{c_w}.
\]
Passing from \(W\) to \(Q_v\) in each local target makes the
coefficient map on \(H^2\) surjective, since local
\(\ell\)-cohomological dimension is two. The kernel dimension
increases by at most the summed drop in the local dimensions.
The last inequality therefore also holds with the new kernel
dimension \(\kappa_2\) and \(h^2(D_v,Q_v)\).

The cone sequence gives
\begin{align*}
 \dim H_C^2-\dim H_C^1+\dim H_C^0
   &=\kappa_2-h^1(G,W)+h^0(G,W)\\
   &\quad+\sum_{v\mid\ell}\bigl(h^1(D_v,Q_v)-h^0(D_v,Q_v)\bigr).
\end{align*}
By the local Euler formula the summand is
\(h^2(D_v,Q_v)+d_v\dim Q_v\). Use the preceding global
inequality and \(\sum_{v\mid\ell}d_v=d\) to obtain
\eqref{eq:fitted-relative-euler}. This gives the relative
relation count on the same local groups and weights that
give the marked first cohomology bounds.
\cite{milne-duality}\label{cor:ordinary-fitted-arithmetic-4}

\end{proof}

\paragraph{The three auxiliary level actions.}
It remains to realize the auxiliary level actions with the prescribed local
characters. The global separating element supplies, by Chebotarev on
the finite quotient recording \(L,P,\bar\epsilon\), degree-one
split places \(v_0\) of \(L^+\) with \(P=1\) on the whole
decomposition group and with nontrivial mod-\(\ell\) cyclotomic
value at Frobenius. Indeed the characters are unramified at
these primes, and the chosen Frobenius fixes \(L\) itself.
Avoid the primes of \(2M\operatorname{disc}(L)\) and any
previous choice. A fixed finite scalar twist keeps all
residual readings on the common scalar line there. If the common residual scalar is \(a\ne0\), the scalar after the cyclotomic twist is different at the chosen Frobenius, because
\[
 \bar\epsilon(\mathrm{Frob}_{v_0})\ne1\quad\Longrightarrow\quad a\bar\epsilon(\mathrm{Frob}_{v_0})\ne a.
\] The same choice makes the local compact-group
condition at \(v_0\) neat, as in
Paragraph~\ref{par:ordinary-functions}.

Use the appropriate scalar twist for each of the rank-two
and two rank-three levels in Section~\ref{sec:linked-seeds}.
At a place above the tame cutoff the inertia type consists
of the diagonal Teichm\"uller characters after this scalar
twist. The same type and the same ordered \(\ell\)-adic
characters are retained at every auxiliary depth of that
level. The two rank-three levels have different
distinguished places. At each of them, the other
distinguished place remains outside the ramification set
of that level. The two actions are compared separately at
their own hyperspecial starting modules before contraction
to the common frame.

\begin{proposition}[Auxiliary actions at the three actual levels]
\label{prop:fitted-three-level-actions}
Under the fitted level conditions, let \(M^{(3)},M^{(0)},M^{(2)}\) denote the modules for the
rank-three primary, rank-three auxiliary and rank-two actions, respectively.
The annihilators for these level actions, computed at empty auxiliary set, are
\[
  \operatorname{Ann}(M^{(3)})=J_3,\qquad
  \operatorname{Ann}(M^{(0)})=J_0,\qquad
  \operatorname{Ann}(M^{(2)})=J_2 .
\] These actions provide the fixed spaces
in the patching argument. The descended containment on the original
host acquires no factor from the fitted degree or number of
Taylor--Wiles places.
\end{proposition}

\begin{proof}
\proofstep{The finite groups at the actual Hecke level.}
Here is the identification at a single level. Enlarge the
allowed ramification set by the selected Taylor--Wiles
places. If \(S\) is the original set of allowed places and \(Q\) the auxiliary set, the quotient at the level in question is
\[
 G_{S\cup Q}\big/\overline{\langle\!\langle I_v:v\in Q\rangle\!\rangle}
      \simeq G_S .
\]
Here the denominator is the closed normal subgroup generated by the images of the added inertia groups.

At hyperspecial compact level at all
new places, the function space and ordinary operator
are exactly the original ones. The ordinary residual blocks
also agree. Indeed the existing operators and those retained
at common spherical places agree, and a classical point
belongs to the residual block exactly when its residual
Frobenius polynomials equal the fixed diagonal polynomials.
Frobenius elements at the common spherical places are dense
up to conjugacy in the quotient. The continuous residual determinant is determined on the quotient by those conjugacy classes. Indeed two such determinants with identical characteristic polynomials on the dense Frobenius set have
\[
 D_1(1-Xg)=D_2(1-Xg)\qquad(g\text{ in the quotient})
\]
by continuity of each coefficient. The local weight character is fixed. Separating classical weights
and finite freeness give the same equality on the integral
ordinary modules. This identifies the empty-auxiliary
annihilator with the Hecke ideal of that level.

At the varying auxiliary levels take the directed union
of the modules of Paragraph~\ref{par:ordinary-functions}.
Right translation on this union is the smooth action
of the product of the auxiliary general linear groups.
The universal determinant and weight ring acts compatibly
on the fixed spaces. For auxiliary levels \(K'\subset K\), smaller level subgroup means a larger fixed space. On the underlying module \(M\) the transition is the inclusion
\[
 i_{K,K'}:M^K\hookrightarrow M^{K'},\qquad v\longmapsto v .
\] For an element \(d\) of the determinant and weight ring form
\[
 A_d=i_{K,K'}\,d_K-d_{K'}\,i_{K,K'} .
\]
At every separating classical specialization, the Frobenius polynomials and the ordered weight characters identify the two determinant actions, so \(A_d\) specializes to zero. On every separating classical specialization \(x\) the defect is zero.

Since these specializations separate the projected ordinary module, an element of the image of this defect has all its specializations zero, so
\[
 \operatorname{im} A_d\subseteq\bigcap_x\ker(\text{specialization at }x)=0 .
\] The same comparison after a change of level by conjugation intertwines the determinant action with translation. Thus these local actions and the
determinant action commute on the directed union.
At the tame level the same modules are projective over
the group ring of the finite \(\ell\)-group quotient by
the argument for the neat level. The augmentation trace
and the two maps with product of split-root factors in
Paragraph~\ref{par:ordinary-level} apply to these fixed
spaces. The integer controlling the local nilpotent relations
depends only on \(n\), and no auxiliary depth or factor is
charged to the final host trace product.

\proofstep{Keeping the original coefficient field in the comparison.}
In the rank-three linked comparison, the primary level
uses the projected rank-three multiplicity computation.
The auxiliary rank-three level and the rank-two level use
the same ordinary construction without the separate
rank-three multiplicity computation. The resulting annihilator ideals, at the three corresponding levels in order, are
\[
 J_3\quad\text{(primary rank three)},\qquad J_0\quad\text{(auxiliary rank three)},\qquad J_2\quad\text{(rank two)} .
\] The patching argument of Section~\ref{sec:marked-patching} on these auxiliary fixed spaces gives ideal containments on the indicated curves. The valuation bound after descent to the
rational host is independent of the dimension bounds on the
fitted field and the number of Taylor--Wiles places.
\label{cor:ordinary-fitted-arithmetic-5}
\end{proof}


\paragraph{A characteristic polynomial on every test curve.}
To control points of the degree-zero algebra we first examine the characteristic
polynomial along a curve in characteristic \(\ell\). Let \(A\) be a complete DVR of
characteristic \(\ell\) with finite residue field and a local map from \(\mathcal O\) through
its residue field \(k\). Take a marked ordinary representation
over \(A\) reducing in the given basis to the diagonal model
\((P,1,P^{-1})\) on \(\ker\eta\). On \(G_{L^+}\) it factors
through the marked quotient. Indeed the killed inertia elements
away from \(\ell\) have \(\eta=P=1\), and the inertia condition
on the point kills these elements. Continuity kills the closed
normal subgroup they generate. Above \(\ell\) the ordered
inertial characters on the frame are \(\Psi^{w_i}\), now
viewed in \(A\), where \(w=(1,0,-1)\). They are one on the
kernel of \(P\) in the marked inertia.

\begin{proposition}[Characteristic polynomials on marked ordinary curves]\label{prop:marked-residual-polynomial}
Let \(A\) be a complete DVR of characteristic \(\ell\) with
finite residue field, with \(\mathcal O\to A\) factoring through
\(k\). Take a marked ordinary curve over \(A\) with residual diagonal
\((P,1,P^{-1})\) on \(\ker\eta\) and ordered ordinary characters
\((\Psi,1,\Psi^{-1})\) on the marked inertia. On the entire quadratic kernel the identity \eqref{eq:tag-h3} reads
\[
 \det(X-r(g))=(X-P(g))(X-1)(X-P(g)^{-1})\qquad(g\in\ker\eta).
\] In the
fixed marked basis the curve is upper or lower triangular,
and its diagonal on the positive component is prescribed
by the residual characters.
\end{proposition}

\begin{proof}
\proofstep{Applying the rank-three support result to the residual curve.}
On the rank-three frame let \(J\) be the image of the Hecke
ideal and \(I_E\) the ideal pulled back from the outer corner
on the exceptional locus. The intersection \(I=J\cap I_E\)
lies below every prime off the two triangular loci. This follows
from propagation on the type-I locus and the rank-two comparison
on the exceptional locus. Write \(\lambda\) for the ideal
generated by the maximal ideal of the coefficients and the
weight differences \(\psi_{vi}(s)-1\) on inertia with \(P(s)=1\).
Both \(J+\lambda\) and \(I_E+\lambda\) have radical the
degree-zero maximal ideal.

For a prime \(\mathfrak p\) the inclusion \(JI_E\subseteq I\) implies
\[
 I+\lambda\subseteq\mathfrak p\quad\Longrightarrow\quad J+\lambda\subseteq\mathfrak p\ \text{ or }\ I_E+\lambda\subseteq\mathfrak p .
\]
Both force the same degree-zero maximal ideal.
The curve kills \(\lambda\). If its kernel is off both triangular
loci, its degree-zero evaluations consequently equal their
residual constants. In the expansion of a characteristic
polynomial each cycle of off-diagonal entries has total weight
zero and vanishes at the residual point. Only products of the
prescribed diagonal entries survive.

\proofstep{The split characteristic polynomial on the fixed subgroup.}
On a triangular curve the diagonal entries are characters.
Divide each by the corresponding residual character. The
quotient takes values in \(1+\mathfrak m_A\). On marked inertia with \(P=1\), any character value \(u\) of the quotient satisfies its ordinary polynomial. Since the coefficient ring here is a domain,
\[
 (u-1)^3=0\quad\Longrightarrow\quad u=1 .
\]
This kills the local abelian generators and, by polarization, their conjugates. At each truncation it is a
character into a finite \(\ell\)-group, and the abelian
class-field condition of the frame gives a single nonzero
exponent \(h\) for all these quotients. Passing through
all truncations gives \(z^h=1\) in \(1+\mathfrak m_A\).

In an equal-characteristic DVR that group is torsion-free.
The prime-to-\(\ell\) part follows by factoring a difference
of powers and the \(\ell\)-part by Frobenius. Thus the three
diagonal characters are fixed in this case as well. If the
comparison was made with enlarged residue coefficients, extend
the test DVR by the corresponding finite residue field
extension, apply the result there and use injectivity on
coefficients. Pulling the polynomial back along the map from
\(G_L\) to the positive component of the frame gives
\[
       \det(X-r(g))=(X-P(g))(X-1)(X-P(g)^{-1})
                        \quad(g\in G_L).
       \tag{H3}\label{eq:tag-h3}
\]
Here and below \(r(g)\) denotes the matrix part on the positive
component. The subgroup \(G_L\) in \eqref{eq:tag-h3} is independent of \(A,r\).

\proofstep{Extend the polynomial identity to the original quadratic kernel.}
For \(g\in\ker\eta\), take \(d\ge1\) with \(g^d\in G_L\).
Let \(\lambda_i\) be the eigenvalues of \(T=r(g)\) in an
algebraic closure of the fraction field, counted with multiplicity.
Triangularizing \(T\) gives the factorization
\[
 \det(X^d-T^d)
   =\prod_i (X^d-\lambda_i^d)
   =\det(X-T)
       \prod_i\left(\sum_{a=0}^{d-1}X^{d-1-a}\lambda_i^a\right).
\]
The first polynomial belongs to \(k[X]\) by \eqref{eq:tag-h3} on \(g^d\),
so all \(\lambda_i\), and hence the coefficients of \(\det(X-T)\),
are algebraic over \(k\). This factorization needs no separability
of \(X^d-\lambda_i^d\). An algebraic element of the domain \(A\)
congruent to \(a\in k\) must equal \(a\), since the field it
generates over \(k\) embeds into the residue field under
reduction. This proves \eqref{eq:tag-h3} on all of \(\ker\eta\), even when
\(\ell\mid d\).

\proofstep{Triangularity in the marked basis.}
We need triangularity in the \emph{marked basis}. Put
\(g^*=cg^{-1}c\). The marked difference in the frame is the diagonal matrix
\[
 r(g_*)-r(g_*^*)=\operatorname{diag}(z,0,-z).
\]
The reduction of \(z\) is \(P(g_*)-P(g_*)^{-1}\ne0\).
Work over the fraction field. An invariant subspace is stable under the marked difference \(D\). The three eigenvalues are distinct, so their coordinate projectors are polynomials in \(D\):
\[
 e_\lambda=\prod_{\nu\ne\lambda}\frac{D-\nu I}{\lambda-\nu},\qquad \lambda,\nu\in\{z,0,-z\}.
\]

Let \(L_1,L_2,L_3\) be the marked coordinate lines, in that order.
A stable subspace is a sum of its images under the three projectors.
Polarization makes its orthogonal complement for the reversed pairing
stable as well, and the complementary planes of the coordinate lines are
\[
 L_1^\perp=L_1\oplus L_2,\qquad L_2^\perp=L_1\oplus L_3,\qquad
 L_3^\perp=L_2\oplus L_3 .
\]
Thus a stable outer line gives a full flag. A stable middle line or
outer plane gives the decomposition \(L_2\oplus(L_1\oplus L_3)\).
If the outer plane is reducible, its two coordinate projectors again
give a stable outer line and a flag. The other possibilities are an irreducible triple
and a block decomposition with irreducible outer pair.

In
the former case Brauer--Nesbitt applied to \eqref{eq:tag-h3} would identify
the irreducible triple with a sum of characters. In the
latter case the middle diagonal entry is an eigenvalue,
algebraic over \(k\) and congruent to one, so it is one.
Since the middle eigenvalue is one, divide the characteristic polynomial in \eqref{eq:tag-h3} by its common factor:
\[
 P_{\rm outer}(X)=\frac{P_{\rm triple}(X)}{X-1}.
\]
The resulting character polynomials identify the semisimplification of the outer pair by Brauer--Nesbitt. Thus every curve is upper or lower triangular.
The equalities of entries descend to \(A\) by injectivity.
Each diagonal entry is then an eigenvalue congruent to its
prescribed value, so equals it on \(\ker\eta\). Multiplication
by \(c\), whose matrix part is one, extends the diagonal
identity to the other coset.

\end{proof}

\paragraph{Finiteness and an integral point in the same orbit.}
The identities on curves can now be used on the original Galois quotient. For the
auxiliary field over \(\mathbb Q(i)\) put \(\Gamma=G_{\Q,S}\), the quotient
by the closed normal subgroup generated by inertia away
from the specified finite set. The surjection
\(G_\Q\to\Gamma\) identifies the image of \(G_K\) with
\(\ker\eta_\Gamma\). The preceding triangularity and diagonal
identities therefore hold on the host. In the criterion for a triangular point the diagonal identity is also
obtained by dividing diagonal characters on \(G_K\).
The inertia conditions make their quotients unramified at
each truncation, and \(\Z[i]\) has class number one.

\begin{construction}
\[
\begin{array}{ll}
\text{monomial on the curve}&\text{reason it vanishes}\\ \hline
\text{contains a weight-zero variable}&\text{the diagonal is the fixed residue diagonal}\\
\text{uses only nonzero weights}&\text{both signs occur to sum to zero}\\
 &\text{one strict sign vanishes by triangularity}
\end{array}
\]
\caption{The closed-fibre test for a nonconstant monomial of weight zero at the Gaussian host. Triangularity is available on each marked residual curve.}\label{fig:marked-host-zero-invariants}
\end{construction}

\begin{proposition}[The integral invariant algebra]\label{prop:marked-invariant-algebra}
Let \(R\) be the complete normalized graded ring of the
Gaussian host just described, with the triangularity and
diagonal identities on all its residual test curves.
The degree-zero subalgebra \(R_0\) is a finite \(\mathcal O\)-module.
For each characteristic-zero point extending the coefficient embedding,
let \(a_i\) be its values on homogeneous generators of the maximal ideal,
of degrees \(d_i\). There exist a finite extension \(O'\) of the coefficient DVR and a grading-torus scalar
\(t\in\operatorname{Frac}(O')^\times\) for which
\[
                  t^{d_i}a_i\in\mathfrak m_{O'}\qquad(\text{every }i).
\]
\end{proposition}

\begin{proof}
\proofstep{Finiteness by testing curves in the closed fibre.}
In an adapted power series
presentation over \(\mathcal O\) take variables \(z_i\) from
entries of \(r-r_0\) on the positive component, with their
entry weights \(w_i\), where \(r_0\) is the split point. They
belong to the maximal ideal. On every residual test curve
the zero-weight variables vanish, as do all variables of one
strict sign. Indeed, if a nonconstant monomial of weight zero has a factor of weight zero it vanishes on the curve. Otherwise its exponents \(a_i\) satisfy
\[
 \sum_i a_iw_i=0,\qquad \sum_i a_i>0 ,
\]
so at least one variable of each strict sign occurs. The variables of one of those signs all vanish on the curve.

Write \(\varpi\) for the coefficient uniformizer. Suppose that a prime ideal
\(\mathfrak p\) contains \(\varpi R\) but does not contain one of these
monomials \(b\). In the complete local domain \(R/\mathfrak p\) we have
\[
 \bar b=b+\mathfrak p\ne0 .
\]
We claim that a further prime
quotient of dimension one can be chosen in which the image of \(b\)
is still nonzero.

This follows by choosing parameters
with \(b\) as the last parameter in the domain, cutting by
the others and taking a minimal prime avoiding \(b\).
Normalization of the resulting complete one-dimensional
local domain gives a complete DVR with finite residue field.
This would be a test curve with \(b\ne0\), contrary to the vanishing on every test curve. Thus some power \(b^m\) belongs to \(\varpi R\) in the ambient graded ring. Extract its degree-zero component:
\[
 b^m\in\varpi R\quad\Longrightarrow\quad b^m\in(\varpi R)_0=\varpi R_0 .
\] The degree-zero monomials topologically generate \(R_0\), since their semigroup is finitely generated. The radical test on these generators thus gives
\[
 \sqrt{\varpi R_0}=\mathfrak m_{R_0}.
\] Noetherianity
gives a power containment. The special fiber has finite
length and the two adic topologies agree. Lifting a basis
of that fiber successively through powers of \(\varpi\)
shows that \(R_0\) is finite over \(\mathcal O\).

\proofstep{An interval of allowable rescalings.}
The characteristic-zero point is initially an evaluation of
the algebra spanned by homogeneous pieces. Write \(a_i\)
for the values of the \(z_i\), with valuation \(v\). Finiteness
of the local algebra \(R_0\) implies \(v>0\) for the values
of its maximal ideal, with \(v(0)=+\infty\). Thus \(v(a_i)>0\)
for weight zero. For \(w_i>0>w_j\), the monomial
\(z_i^{|w_j|}z_j^{w_i}\) has strictly positive valuation.
Consequently the finitely many strict constraints
\[
             -v(a_i)/w_i < b < v(a_j)/|w_j|
                \qquad(w_i>0>w_j,\ a_ia_j\ne0)
\]
have a common rational solution, also when only one sign
occurs by keeping the corresponding one-sided constraints.
Choose \(t\in\overline{\Q}_\ell^\times\) with \(v(t)=b\).
All the \(t^{w_i}a_i\) lie in the maximal ideal of one finite
coefficient DVR. Evaluation of power series converges there.

\proofstep{Evaluation of homogeneous series.}
To check agreement with the given point on a homogeneous
series, decompose it into finitely many monomials of its weight
times series of weight zero. Such generators exist by the
coordinatewise minimal exponent vectors and Dickson's lemma.
Evaluation is continuous on \(R_0\) by finiteness and the
adic power containment. The two evaluations on weight \(j\)
therefore differ exactly by \(t^j\). In matrix terms the
integral representation and the given point are conjugate
by \(\operatorname{diag}(t,1,t^{-1})\), with the semidirect
conjugation law on the negative coset. Both use the image of
the \emph{same} \(g_*\).
\end{proof}
